% === ОБЩИЕ НАСТРОЙКИ ДЛЯ ВСЕХ ШАБЛОНОВ ===

\newcommand{\docTitle}{Документ}
\newcommand{\docAuthor}{Организация}
\newcommand{\docHeader}{Заголовок документа}
\newcommand{\docObject}{}
\newcommand{\docCategory}{}
\newcommand{\docCity}{г. Омск}
\newcommand{\docDate}{\today}
\newcommand{\docOrganization}{Организация}
\newcommand{\docType}{Документ}
\newcommand{\docSubtitle}{}
\newcommand{\docSubject}{Тема документа}
\newcommand{\docKeywords}{}
\newcommand{\docYear}{\the\year}

% 2. Шрифты и языки (XeLaTeX)
\usepackage{fontspec}
\usepackage{polyglossia}
\setmainlanguage{russian}
\setotherlanguage{english}

\defaultfontfeatures{Ligatures=TeX, Script=Cyrillic, Path=C:/Windows/Fonts/}
\setmainfont{times.ttf}[
  UprightFont = *,
  BoldFont = timesbd.ttf,
  ItalicFont = timesi.ttf,
  BoldItalicFont = timesbi.ttf
]
\setsansfont{arial.ttf}[
  UprightFont = *,
  BoldFont = arialbd.ttf,
  ItalicFont = ariali.ttf,
  BoldItalicFont = arialbi.ttf
]
\setmonofont{cour.ttf}[
  UprightFont = *,
  BoldFont = courbd.ttf,
  ItalicFont = couri.ttf,
  BoldItalicFont = courbi.ttf
]

% 3. Основные пакеты
\usepackage{amsmath, amssymb, amsfonts}
\usepackage{graphicx}
\usepackage{float}
\usepackage{array}
\usepackage{tabularx}
\usepackage{longtable}
\usepackage{booktabs}
\usepackage{multirow}

% 4. Списки
\usepackage{enumitem}
\setlist[itemize]{label={---}, leftmargin=*, topsep=4pt, itemsep=2pt}
\setlist[enumerate]{label={\arabic*.}, leftmargin=*, topsep=4pt, itemsep=2pt}

% 5. Ссылки и PDF-метаданные (теперь \docTitle уже определён!)
\usepackage{hyperref}
\hypersetup{
    colorlinks=false,
    pdfborder={0 0 0},
    unicode=true,
    pdftitle={\docTitle},
    pdfauthor={\docAuthor},
    pdfsubject={\docSubject}
}

% 6. Совместимость с Pandoc
\providecommand{\tightlist}{%
  \setlength{\itemsep}{0pt}\setlength{\parskip}{0pt}}

% 7. Абзацы
\setlength{\parindent}{1.25cm}
\setlength{\parskip}{0pt}

% 8. Команды для ссылок на нормативку
\newcommand{\gost}[1]{ГОСТ\,#1}
\newcommand{\fstek}[1]{ФСТЭК\,#1}
\newcommand{\fz}[1]{ФЗ\,#1}

% 9. Термины
\newcommand{\term}[1]{\textbf{#1}}
\newcommand{\abbr}[2]{#1 (#2)}

% 10. Стиль для основной страницы
\newcommand{\mainpagestyle}{%
    \pagestyle{fancy}%
    \fancyhf{}%
    \fancyhead[C]{\small \docHeader}%
    \fancyfoot[C]{\thepage}%
    \renewcommand{\headrulewidth}{0.5pt}%
    \renewcommand{\footrulewidth}{0pt}%
}